\documentclass[11pt,a4paper]{article}
\usepackage[utf8]{inputenc}
\usepackage[french]{babel}
\usepackage[T1]{fontenc}
\usepackage{geometry}
\geometry{left=2.5cm,right=2.5cm,top=2.5cm,bottom=2.5cm}
\usepackage{tikz}
\usetikzlibrary{shapes.geometric, arrows.meta, positioning, fit, backgrounds, shadows, decorations.pathreplacing, calc}
\usepackage{xcolor}
\usepackage{tcolorbox}
\tcbuselibrary{skins,breakable}
\usepackage{enumitem}
\usepackage{tabularx}
\usepackage{booktabs}
\usepackage{multirow}
\usepackage{graphicx}
\usepackage{amsmath}
\usepackage{amssymb}

% Définition des couleurs pour chaque zone
\definecolor{zone1color}{RGB}{66, 135, 245}    % Bleu - Data Security
\definecolor{zone2color}{RGB}{76, 175, 80}     % Vert - Training Security
\definecolor{zone3color}{RGB}{255, 193, 7}     % Jaune - Validation
\definecolor{zone4color}{RGB}{255, 87, 34}     % Orange - Production
\definecolor{zone5color}{RGB}{156, 39, 176}    % Violet - Governance

\title{\Huge\bfseries Architecture de Sécurisation à 5 Zones\\[0.5em]
\Large Pour Systèmes d'Intelligence Artificielle Critiques}
\author{Méthode de sécurisation appliquée au maintien de l'ordre}
\date{}

\begin{document}

\maketitle

\begin{abstract}
Ce document présente l'architecture complète de sécurisation à 5 zones développée pour garantir la robustesse, la confidentialité et la fiabilité des systèmes d'intelligence artificielle déployés dans des contextes critiques. L'architecture couvre l'ensemble du cycle de vie d'un modèle d'IA, de la collecte des données jusqu'à la gouvernance post-déploiement, en intégrant des mécanismes de défense contre les attaques adversariales, le data poisoning, et les violations de confidentialité.
\end{abstract}

\tableofcontents
\newpage

\section{Vue d'ensemble de l'architecture}

\subsection{Principe de défense en profondeur}

L'architecture proposée s'appuie sur le principe de \textbf{défense en profondeur} (Defense in Depth), où chaque zone constitue une barrière de sécurité indépendante. Cette approche multicouche garantit qu'une compromission à un niveau n'entraîne pas l'effondrement total du système de sécurité.

\begin{figure}[h!]
\centering
\begin{tikzpicture}[
    scale=0.9,
    every node/.style={font=\small},
    zone/.style={
        rectangle,
        rounded corners=3mm,
        draw=black,
        very thick,
        minimum width=12cm,
        minimum height=2cm,
        text width=11cm,
        align=left,
        drop shadow
    },
    arrow/.style={
        -Stealth,
        very thick,
        draw=black!70
    }
]

% Zone 1 - Data Security
\node[zone, fill=zone1color!20] (z1) at (0,0) {
    \textbf{\Large Zone 1: Data Security}\\[0.2em]
    \textbullet~Validation et vérification des données\\
    \textbullet~Détection de data poisoning\\
    \textbullet~Chiffrement et anonymisation\\
    \textbullet~Contrôle d'intégrité (hashing)
};

% Zone 2 - Training Security
\node[zone, fill=zone2color!20, below=0.5cm of z1] (z2) {
    \textbf{\Large Zone 2: Training Security}\\[0.2em]
    \textbullet~Adversarial Training (TRADES)\\
    \textbullet~Differential Privacy (ε,δ)-DP\\
    \textbullet~Regularisation robuste\\
    \textbullet~Isolation environnement d'entraînement
};

% Zone 3 - Validation
\node[zone, fill=zone3color!20, below=0.5cm of z2] (z3) {
    \textbf{\Large Zone 3: Validation \& Testing}\\[0.2em]
    \textbullet~Tests adversariaux (FGSM, PGD, AutoAttack)\\
    \textbullet~Audit de biais et équité\\
    \textbullet~Certification de robustesse\\
    \textbullet~Tests d'intégration sécurisés
};

% Zone 4 - Production
\node[zone, fill=zone4color!20, below=0.5cm of z3] (z4) {
    \textbf{\Large Zone 4: Production Deployment}\\[0.2em]
    \textbullet~Détection d'attaques en temps réel\\
    \textbullet~Input filtering et sanitization\\
    \textbullet~Rate limiting et authentification\\
    \textbullet~Containérisation et isolation (Docker)
};

% Zone 5 - Governance
\node[zone, fill=zone5color!20, below=0.5cm of z4] (z5) {
    \textbf{\Large Zone 5: Governance \& Monitoring}\\[0.2em]
    \textbullet~Monitoring continu (Prometheus/Grafana)\\
    \textbullet~Audit logging immutable\\
    \textbullet~Alerting et incident response\\
    \textbullet~Conformité réglementaire (RGPD)
};

% Flèches entre zones
\draw[arrow] (z1.south) -- (z2.north);
\draw[arrow] (z2.south) -- (z3.north);
\draw[arrow] (z3.south) -- (z4.north);
\draw[arrow] (z4.south) -- (z5.north);

% Feedback loops
\draw[arrow, dashed, red, bend right=60] (z5.west) to node[left, text=red] {Feedback} (z1.west);
\draw[arrow, dashed, red, bend right=50] (z3.east) to node[right, text=red] {Retraining} (z2.east);

\end{tikzpicture}
\caption{Vue d'ensemble de l'architecture à 5 zones de sécurité}
\label{fig:overview}
\end{figure}

\subsection{Principes directeurs}

\begin{tcolorbox}[colback=blue!5, colframe=blue!50!black, title=\textbf{Principes Fondamentaux}]
\begin{enumerate}[leftmargin=*]
    \item \textbf{Isolation}: Chaque zone est isolée avec des frontières de sécurité claires
    \item \textbf{Validation}: Validation stricte des entrées/sorties à chaque transition
    \item \textbf{Traçabilité}: Logging complet et audit trail pour toutes les opérations
    \item \textbf{Réversibilité}: Capacité de rollback à tout niveau en cas de compromission
    \item \textbf{Monitoring}: Surveillance continue avec détection d'anomalies
\end{enumerate}
\end{tcolorbox}

\newpage

\section{Zone 1: Data Security}

\subsection{Objectif et périmètre}

La Zone 1 constitue le premier rempart de sécurité, garantissant que seules des données propres, validées et anonymisées entrent dans le pipeline d'entraînement.

\begin{figure}[h!]
\centering
\begin{tikzpicture}[
    node distance=1.5cm,
    box/.style={rectangle, draw, rounded corners, fill=blue!10, minimum width=3cm, minimum height=1cm, align=center, drop shadow},
    process/.style={rectangle, draw, fill=green!10, minimum width=2.5cm, minimum height=0.8cm, align=center},
    decision/.style={diamond, draw, fill=yellow!10, aspect=2, align=center},
    arrow/.style={-Stealth, thick}
]

% Input
\node[box, fill=gray!20] (input) {Données brutes\\(Images, Métadonnées)};

% Data Validation Pipeline
\node[process, below=of input] (hash) {Vérification\\Intégrité (SHA-256)};
\node[process, below=of hash] (poison) {Détection\\Poisoning};
\node[process, below=of poison] (privacy) {Anonymisation\\EXIF removal};
\node[process, below=of privacy] (encrypt) {Chiffrement\\AES-256};

% Decision
\node[decision, right=3cm of poison] (valid) {Données\\Valides?};

% Output
\node[box, fill=green!20, below=of encrypt] (clean) {Données propres\\et sécurisées};
\node[box, fill=red!20, right=of valid] (reject) {Rejet\\+ Alerte};

% Arrows
\draw[arrow] (input) -- (hash);
\draw[arrow] (hash) -- (poison);
\draw[arrow] (poison) -- (valid);
\draw[arrow] (valid) -- node[above] {Oui} (privacy);
\draw[arrow] (privacy) -- (encrypt);
\draw[arrow] (encrypt) -- (clean);
\draw[arrow] (valid) -- node[above] {Non} (reject);

% Statistical Tests Box
\node[box, fill=orange!10, left=3cm of poison, text width=2.5cm, font=\footnotesize] (tests) {
    \textbf{Tests Statistiques:}\\
    $\bullet$ Chi-Square\\
    $\bullet$ Z-score\\
    $\bullet$ Kolmogorov-Smirnov\\
    $\bullet$ Isolation Forest
};
\draw[arrow, dashed] (tests) -- (poison);

\end{tikzpicture}
\caption{Pipeline de validation des données - Zone 1}
\label{fig:zone1_pipeline}
\end{figure}

\subsection{Mécanismes de sécurité}

\begin{table}[h!]
\centering
\caption{Mécanismes de défense de la Zone 1}
\label{tab:zone1_defenses}
\begin{tabularx}{\textwidth}{|l|X|X|}
\hline
\textbf{Mécanisme} & \textbf{Technologie} & \textbf{Protection} \\ \hline
Vérification d'intégrité & Hash SHA-256 & Détection de corruption/modification \\ \hline
Détection poisoning & Isolation Forest (scikit-learn) & Identification d'échantillons malveillants \\ \hline
Anonymisation & Suppression métadonnées EXIF & Protection vie privée \\ \hline
Chiffrement stockage & AES-256 (Fernet) & Confidentialité at-rest \\ \hline
Tests statistiques & Chi-Square, KS-test, Z-score & Détection d'anomalies distribution \\ \hline
\end{tabularx}
\end{table}

\subsection{Implémentation}

\textbf{Fichier}: \texttt{src/data/data\_verifier.py}

\begin{tcolorbox}[colback=gray!5, colframe=gray!50, title=Extrait de code - DataVerifier]
\begin{verbatim}
class DataVerifier:
    def verify_dataset(self, dataset_path):
        # 1. Collecter statistiques
        statistics = self._collect_statistics(dataset_path)

        # 2. Tests statistiques
        statistical_tests = self._run_statistical_tests(
            dataset_path, statistics
        )

        # 3. Détecter anomalies
        anomalies = self._detect_anomalies(
            statistics, statistical_tests
        )

        # 4. Score de qualité
        quality_score = self._calculate_quality_score(
            statistics, statistical_tests, anomalies
        )

        return VerificationReport(...)
\end{verbatim}
\end{tcolorbox}

\newpage

\section{Zone 2: Training Security}

\subsection{Objectif et périmètre}

La Zone 2 sécurise le processus d'entraînement en intégrant des défenses adversariales directement dans l'objectif d'optimisation et en garantissant la confidentialité des données d'entraînement.

\begin{figure}[h!]
\centering
\begin{tikzpicture}[
    node distance=1.5cm,
    box/.style={rectangle, draw, rounded corners, fill=green!10, minimum width=3.5cm, minimum height=1cm, align=center, drop shadow},
    data/.style={cylinder, draw, shape border rotate=90, fill=blue!10, minimum width=2cm, minimum height=1.5cm, align=center},
    arrow/.style={-Stealth, thick}
]

% Data
\node[data] (train_data) {Dataset\\Propre\\(Zone 1)};

% Training Components
\node[box, right=2cm of train_data] (model) {Modèle\\ResNet50};
\node[box, above=1cm of model] (at) {Adversarial\\Training\\(TRADES)};
\node[box, below=1cm of model] (dp) {Differential\\Privacy\\(ε,δ)-DP};
\node[box, right=2cm of model] (robust) {Modèle\\Robuste\\Entraîné};

% Loss computation
\node[box, fill=yellow!10, above right=0.5cm and 0.5cm of model, text width=3cm] (loss) {
    \textbf{TRADES Loss:}\\
    $L = L_{CE}(x) +$\\
    $\beta \cdot L_{KL}(x, x_{adv})$
};

% DP mechanism
\node[box, fill=orange!10, below right=0.5cm and 0.5cm of model, text width=3cm] (noise) {
    \textbf{DP Noise:}\\
    $\nabla \leftarrow \nabla + \mathcal{N}(0, \sigma^2)$\\
    Clip: $||\nabla|| \leq C$
};

% Arrows
\draw[arrow] (train_data) -- (model);
\draw[arrow] (at) -- (model);
\draw[arrow] (dp) -- (model);
\draw[arrow] (model) -- (robust);
\draw[arrow, dashed] (loss) -- (model);
\draw[arrow, dashed] (noise) -- (model);

% Adversarial examples generation
\node[box, fill=red!10, above=0.5cm of at, text width=3cm, font=\footnotesize] (pgd) {
    \textbf{PGD Attack:}\\
    $x_{adv}^{t+1} = \Pi_{\varepsilon}(x_{adv}^t + \alpha \cdot sign(\nabla_x L))$
};
\draw[arrow, dashed] (pgd) -- (at);

\end{tikzpicture}
\caption{Architecture d'entraînement sécurisé - Zone 2}
\label{fig:zone2_training}
\end{figure}

\subsection{Mécanismes de sécurité}

\begin{table}[h!]
\centering
\caption{Mécanismes de défense de la Zone 2}
\label{tab:zone2_defenses}
\begin{tabularx}{\textwidth}{|l|X|X|}
\hline
\textbf{Mécanisme} & \textbf{Technologie} & \textbf{Protection} \\ \hline
Adversarial Training & TRADES (Zhang et al. 2019) & Robustesse contre attaques $L_\infty$ \\ \hline
Differential Privacy & Opacus (PyTorch) & Confidentialité données d'entraînement \\ \hline
Gradient Clipping & $||\nabla|| \leq C$ & Stabilité numérique + DP \\ \hline
KL Divergence & Régularisation logits & Prévention gradient masking \\ \hline
Environnement isolé & Docker container & Isolation système \\ \hline
\end{tabularx}
\end{table}

\subsection{Formulation mathématique}

\begin{tcolorbox}[colback=blue!5, colframe=blue!50!black, title=\textbf{TRADES Loss Function}]
La fonction de perte TRADES combine la précision sur données propres et la robustesse adversariale:

\begin{equation}
\mathcal{L}_{TRADES} = \mathbb{E}_{(x,y)\sim\mathcal{D}} \left[ L_{CE}(f_\theta(x), y) + \beta \cdot D_{KL}(f_\theta(x) || f_\theta(x_{adv})) \right]
\end{equation}

où:
\begin{itemize}
    \item $L_{CE}$: Cross-Entropy Loss sur données propres
    \item $D_{KL}$: Divergence de Kullback-Leibler
    \item $\beta$: Coefficient de trade-off (typiquement $\beta = 6.0$)
    \item $x_{adv}$: Exemple adversarial généré par PGD
\end{itemize}
\end{tcolorbox}

\subsection{Implémentation}

\textbf{Fichier}: \texttt{src/experiments/secured/train\_secured\_trades.py}

\begin{tcolorbox}[colback=gray!5, colframe=gray!50, title=Extrait de code - TRADES Training]
\begin{verbatim}
def _trades_loss(self, logits_clean, logits_adv, labels, beta):
    # Cross-Entropy Loss sur données propres
    ce_loss = self.ce_criterion(logits_clean, labels)

    # KL Divergence Loss entre propres et adversariaux
    prob_clean = F.log_softmax(logits_clean, dim=1)
    prob_adv = F.softmax(logits_adv, dim=1)
    kl_loss = self.kl_criterion(prob_clean, prob_adv)

    # TRADES Loss combinée
    total_loss = ce_loss + beta * kl_loss

    return total_loss, ce_loss, kl_loss
\end{verbatim}
\end{tcolorbox}

\newpage

\section{Zone 3: Validation \& Testing}

\subsection{Objectif et périmètre}

La Zone 3 effectue une validation rigoureuse du modèle entraîné à travers une batterie de tests adversariaux, d'audits de biais, et de certifications de robustesse avant tout déploiement en production.

\begin{figure}[h!]
\centering
\begin{tikzpicture}[
    node distance=1.5cm,
    box/.style={rectangle, draw, rounded corners, fill=yellow!10, minimum width=3cm, minimum height=0.8cm, align=center, drop shadow},
    test/.style={rectangle, draw, fill=orange!10, minimum width=2.5cm, minimum height=0.7cm, align=center},
    arrow/.style={-Stealth, thick}
]

% Input model
\node[box, fill=green!20] (model) {Modèle\\Entraîné\\(Zone 2)};

% Test categories
\node[box, below right=1cm and -1cm of model] (cat1) {Tests\\Adversariaux};
\node[box, below=0.5cm of cat1] (cat2) {Audit\\Équité/Biais};
\node[box, below=0.5cm of cat2] (cat3) {Tests\\Performance};
\node[box, below=0.5cm of cat3] (cat4) {Certification\\Robustesse};

% Specific tests
\node[test, right=2cm of cat1] (fgsm) {FGSM\\$\varepsilon=0.031$};
\node[test, right=0.2cm of fgsm] (pgd) {PGD\\10 iter};
\node[test, right=0.2cm of pgd] (aa) {AutoAttack};

\node[test, right=2cm of cat2] (fairness) {Disparité\\Démographique};
\node[test, right=0.2cm of fairness] (fpr) {FPR/FNR\\par classe};

\node[test, right=2cm of cat3] (acc) {Accuracy};
\node[test, right=0.2cm of acc] (f1) {F1-Score};
\node[test, right=0.2cm of f1] (roc) {ROC-AUC};

\node[test, right=2cm of cat4] (cert) {Score\\Robustesse};
\node[test, right=0.2cm of cert] (grade) {Security\\Grade};

% Decision
\node[box, fill=red!20, below=3cm of cat3] (fail) {\textbf{ÉCHEC}\\Retour Zone 2};
\node[box, fill=green!30, right=2cm of fail] (pass) {\textbf{SUCCÈS}\\Déploiement Zone 4};

% Arrows
\draw[arrow] (model) -- (cat1);
\draw[arrow] (model) -- (cat2);
\draw[arrow] (model) -- (cat3);
\draw[arrow] (model) -- (cat4);

\draw[arrow] (cat1) -- (fgsm);
\draw[arrow] (cat2) -- (fairness);
\draw[arrow] (cat3) -- (acc);
\draw[arrow] (cat4) -- (cert);

\draw[arrow] (cat1) |- (fail);
\draw[arrow] (cat2) |- (fail);
\draw[arrow] (cat3) |- (fail);
\draw[arrow] (cat4) -- (pass);

\end{tikzpicture}
\caption{Pipeline de validation et testing - Zone 3}
\label{fig:zone3_validation}
\end{figure}

\subsection{Métriques de certification}

\begin{table}[h!]
\centering
\caption{Critères de certification de sécurité}
\label{tab:zone3_certification}
\begin{tabularx}{\textwidth}{|l|c|X|}
\hline
\textbf{Critère} & \textbf{Seuil} & \textbf{Description} \\ \hline
Clean Accuracy & $\geq 95\%$ & Performance sur données non altérées \\ \hline
FGSM Robustness & $\geq 80\%$ & Résistance à FGSM ($\varepsilon=0.031$) \\ \hline
PGD Robustness & $\geq 80\%$ & Résistance à PGD (10 itérations) \\ \hline
Disparité FPR & $\leq 5\%$ & Équité entre classes (False Positive Rate) \\ \hline
Disparité FNR & $\leq 5\%$ & Équité entre classes (False Negative Rate) \\ \hline
Latence CPU & $\leq 200ms$ & Performance temps réel \\ \hline
\end{tabularx}
\end{table}

\subsection{Grades de sécurité}

\begin{tcolorbox}[colback=green!5, colframe=green!50!black, title=\textbf{Système de grading}]
\begin{itemize}
    \item \textbf{Grade A+}: Tous critères passés + PGD Robustness $\geq 90\%$ $\rightarrow$ \textcolor{green}{\textbf{Production Ready}}
    \item \textbf{Grade A}: Tous critères passés $\rightarrow$ \textcolor{green}{\textbf{Production Ready}}
    \item \textbf{Grade B}: 2/3 critères passés $\rightarrow$ \textcolor{orange}{\textbf{Review Required}}
    \item \textbf{Grade C}: 1/3 critères passés $\rightarrow$ \textcolor{red}{\textbf{Not Ready}}
    \item \textbf{Grade F}: Aucun critère passé $\rightarrow$ \textcolor{red}{\textbf{Critical Failure}}
\end{itemize}
\end{tcolorbox}

\subsection{Implémentation}

\textbf{Fichier}: \texttt{src/experiments/secured/attack\_secured.py}

\newpage

\section{Zone 4: Production Deployment}

\subsection{Objectif et périmètre}

La Zone 4 déploie le modèle certifié en production avec des défenses actives contre les attaques en temps réel, un monitoring continu, et des mécanismes de résilience.

\begin{figure}[h!]
\centering
\begin{tikzpicture}[
    node distance=1cm,
    box/.style={rectangle, draw, rounded corners, fill=orange!10, minimum width=2.8cm, minimum height=0.9cm, align=center, drop shadow, font=\small},
    component/.style={rectangle, draw, fill=blue!10, minimum width=2.5cm, minimum height=0.7cm, align=center, font=\footnotesize},
    arrow/.style={-Stealth, thick}
]

% External layer
\node[box, fill=red!20] (client) {Client\\Requête};

% Security layer
\node[box, below=1.5cm of client] (waf) {Web Application\\Firewall (WAF)};
\node[component, below=0.3cm of waf] (auth) {Authentification\\JWT};
\node[component, below=0.3cm of auth] (rate) {Rate Limiting\\(100 req/min)};
\node[component, below=0.3cm of rate] (filter) {Input Filtering\\Sanitization};

% API Layer
\node[box, below=0.5cm of filter, fill=green!20] (api) {API FastAPI\\Docker Container};

% Model Layer
\node[box, below=0.5cm of api] (detector) {Adversarial\\Detector};
\node[box, left=1cm of detector] (model) {Modèle\\Robuste};
\node[box, right=1cm of detector] (ensemble) {Ensemble\\Defense};

% Output Layer
\node[box, below=1cm of detector, fill=green!30] (response) {Réponse\\Sécurisée};

% Monitoring
\node[box, fill=purple!20, right=3cm of api, text width=2.5cm] (monitor) {Monitoring\\Prometheus\\Grafana};

% Logging
\node[box, fill=purple!20, right=3cm of response, text width=2.5cm] (logs) {Audit Logging\\Immutable\\Logs};

% Arrows
\draw[arrow] (client) -- (waf);
\draw[arrow] (waf) -- (auth);
\draw[arrow] (auth) -- (rate);
\draw[arrow] (rate) -- (filter);
\draw[arrow] (filter) -- (api);
\draw[arrow] (api) -- (detector);
\draw[arrow] (detector) -- (model);
\draw[arrow] (detector) -- (ensemble);
\draw[arrow] (model) -- (response);
\draw[arrow] (ensemble) -- (response);
\draw[arrow] (response) -- (client);

% Monitoring arrows
\draw[arrow, dashed, purple] (waf) -- (monitor);
\draw[arrow, dashed, purple] (api) -- (monitor);
\draw[arrow, dashed, purple] (detector) -- (monitor);
\draw[arrow, dashed, purple] (model) -- (logs);
\draw[arrow, dashed, purple] (response) -- (logs);

\end{tikzpicture}
\caption{Architecture de déploiement en production - Zone 4}
\label{fig:zone4_deployment}
\end{figure}

\subsection{Mécanismes de sécurité}

\begin{table}[h!]
\centering
\caption{Mécanismes de défense de la Zone 4}
\label{tab:zone4_defenses}
\begin{tabularx}{\textwidth}{|l|X|X|}
\hline
\textbf{Mécanisme} & \textbf{Technologie} & \textbf{Protection} \\ \hline
WAF & Nginx + ModSecurity & Filtrage requêtes malicieuses \\ \hline
Authentification & JWT (JSON Web Tokens) & Contrôle d'accès \\ \hline
Rate Limiting & Redis + Token Bucket & Protection DoS \\ \hline
Input Validation & FastAPI Pydantic & Validation schéma données \\ \hline
Adversarial Detection & Statistical tests + MagNet & Détection attaques temps réel \\ \hline
Containerisation & Docker + security profiles & Isolation et sandboxing \\ \hline
Circuit Breaker & Resilience4j pattern & Résilience en cas de surcharge \\ \hline
\end{tabularx}
\end{table}

\subsection{Architecture API}

\textbf{Fichier}: \texttt{src/api/app.py}

\begin{tcolorbox}[colback=gray!5, colframe=gray!50, title=Extrait de code - FastAPI Secured Endpoint]
\begin{verbatim}
@app.post("/detect/secured")
async def detect_secured(
    file: UploadFile,
    current_user: User = Depends(get_current_user)
):
    # 1. Rate limiting check
    await check_rate_limit(current_user.id)

    # 2. Input validation
    image = await validate_and_preprocess(file)

    # 3. Adversarial detection
    is_adversarial = adversarial_detector.detect(image)
    if is_adversarial:
        raise HTTPException(403, "Adversarial input detected")

    # 4. Model inference
    prediction = secured_model.predict(image)

    # 5. Audit logging
    audit_log.record(user=current_user, prediction=prediction)

    return {"prediction": prediction, "confidence": ...}
\end{verbatim}
\end{tcolorbox}

\newpage

\section{Zone 5: Governance \& Monitoring}

\subsection{Objectif et périmètre}

La Zone 5 assure la gouvernance continue du système déployé à travers un monitoring en temps réel, des audits réguliers, et une conformité réglementaire permanente.

\begin{figure}[h!]
\centering
\begin{tikzpicture}[
    node distance=1.2cm,
    box/.style={rectangle, draw, rounded corners, fill=purple!10, minimum width=3cm, minimum height=1cm, align=center, drop shadow},
    metric/.style={rectangle, draw, fill=blue!10, minimum width=2.5cm, minimum height=0.7cm, align=center, font=\footnotesize},
    arrow/.style={-Stealth, thick}
]

% Central monitoring hub
\node[box, fill=purple!30] (hub) {\textbf{Monitoring Hub}\\Prometheus + Grafana};

% Data sources
\node[metric, above left=1cm and 2cm of hub] (perf) {Métriques\\Performance};
\node[metric, above=1.5cm of hub] (security) {Métriques\\Sécurité};
\node[metric, above right=1cm and 2cm of hub] (quality) {Métriques\\Qualité};

% Detailed metrics
\node[metric, above=0.3cm of perf, text width=2.3cm, font=\scriptsize] (perf_det) {
    Latence (P50/P95)\\
    Throughput\\
    Resource usage
};

\node[metric, above=0.3cm of security, text width=2.3cm, font=\scriptsize] (sec_det) {
    Attaques détectées\\
    Auth failures\\
    Anomalies
};

\node[metric, above=0.3cm of quality, text width=2.3cm, font=\scriptsize] (qual_det) {
    Accuracy glissante\\
    Drift détection\\
    Biais mesurés
};

% Actions
\node[box, below left=1cm and 1.5cm of hub, fill=red!20] (alert) {Alerting\\Email/SMS};
\node[box, below=1.5cm of hub, fill=orange!20] (report) {Reporting\\Automatisé};
\node[box, below right=1cm and 1.5cm of hub, fill=green!20] (comply) {Conformité\\RGPD};

% Audit
\node[box, below=2.5cm of hub, fill=yellow!20, text width=8cm] (audit) {
    \textbf{Audit Trail Immutable}\\
    Blockchain-inspired logging $\cdot$ WORM storage $\cdot$ Forensics
};

% Arrows
\draw[arrow] (perf_det) -- (perf);
\draw[arrow] (sec_det) -- (security);
\draw[arrow] (qual_det) -- (quality);

\draw[arrow] (perf) -- (hub);
\draw[arrow] (security) -- (hub);
\draw[arrow] (quality) -- (hub);

\draw[arrow] (hub) -- (alert);
\draw[arrow] (hub) -- (report);
\draw[arrow] (hub) -- (comply);
\draw[arrow] (hub) -- (audit);

% Feedback loop
\draw[arrow, dashed, red, bend right=45] (alert.west) to node[left, text=red, font=\scriptsize] {Incident\\Response} +(-3,-2);

\end{tikzpicture}
\caption{Architecture de gouvernance et monitoring - Zone 5}
\label{fig:zone5_governance}
\end{figure}

\subsection{Dashboard Grafana}

\begin{table}[h!]
\centering
\caption{Métriques clés du dashboard}
\label{tab:zone5_metrics}
\small
\begin{tabularx}{\textwidth}{|l|l|X|}
\hline
\textbf{Catégorie} & \textbf{Métrique} & \textbf{Seuil d'alerte} \\ \hline
\multirow{3}{*}{Performance}
    & Latence P95 & $> 200ms$ \\
    & Throughput & $< 10$ req/s \\
    & CPU Usage & $> 80\%$ \\ \hline
\multirow{3}{*}{Sécurité}
    & Attaques/min & $> 5$ \\
    & Auth failures & $> 10$/10min \\
    & Adversarial rate & $> 1\%$ \\ \hline
\multirow{3}{*}{Qualité}
    & Accuracy (24h) & $< 85\%$ \\
    & Model drift score & $> 0.1$ \\
    & FPR disparity & $> 5\%$ \\ \hline
\end{tabularx}
\end{table}

\subsection{Conformité réglementaire}

\begin{tcolorbox}[colback=blue!5, colframe=blue!50!black, title=\textbf{Conformité RGPD - Checklist}]
\begin{itemize}[leftmargin=*]
    \item[$\square$] \textbf{Droit à l'information}: Documentation transparente du système
    \item[$\square$] \textbf{Droit d'accès}: API pour consulter les données personnelles
    \item[$\square$] \textbf{Droit à l'effacement}: Procédure de suppression des données
    \item[$\square$] \textbf{Minimisation des données}: Collecte limitée au strict nécessaire
    \item[$\square$] \textbf{Limitation de la conservation}: Purge automatique après 6 mois
    \item[$\square$] \textbf{Intégrité et confidentialité}: Chiffrement AES-256 + TLS 1.3
    \item[$\square$] \textbf{DPIA}: Data Protection Impact Assessment complétée
    \item[$\square$] \textbf{Registre des traitements}: Tenu à jour et auditable
\end{itemize}
\end{tcolorbox}

\newpage

\section{Flux de données inter-zones}

\subsection{Vue d'ensemble du flux}

\begin{figure}[h!]
\centering
\begin{tikzpicture}[
    scale=0.85,
    every node/.style={font=\footnotesize},
    zone/.style={rectangle, draw, rounded corners, minimum width=2.5cm, minimum height=1.8cm, align=center, drop shadow},
    data/.style={ellipse, draw, fill=gray!20, minimum width=1.5cm, minimum height=0.8cm, align=center},
    arrow/.style={-Stealth, very thick}
]

% Zones
\node[zone, fill=zone1color!20] (z1) at (0,0) {\textbf{Zone 1}\\Data\\Security};
\node[zone, fill=zone2color!20] (z2) at (4,0) {\textbf{Zone 2}\\Training\\Security};
\node[zone, fill=zone3color!20] (z3) at (8,0) {\textbf{Zone 3}\\Validation};
\node[zone, fill=zone4color!20] (z4) at (6,-3) {\textbf{Zone 4}\\Production};
\node[zone, fill=zone5color!20] (z5) at (2,-3) {\textbf{Zone 5}\\Governance};

% Data flows
\node[data] (d1) at (2,0) {Données\\Propres};
\node[data] (d2) at (6,0) {Modèle\\Entraîné};
\node[data] (d3) at (9,-1.5) {Modèle\\Certifié};
\node[data] (d4) at (4,-3) {Métriques\\Logs};
\node[data] (d5) at (0,-1.5) {Alertes\\Reports};

% Arrows between zones
\draw[arrow, zone1color] (z1) -- (d1);
\draw[arrow, zone1color] (d1) -- (z2);
\draw[arrow, zone2color] (z2) -- (d2);
\draw[arrow, zone2color] (d2) -- (z3);
\draw[arrow, zone3color] (z3) -- (d3);
\draw[arrow, zone3color] (d3) -- (z4);
\draw[arrow, zone4color] (z4) -- (d4);
\draw[arrow, zone4color] (d4) -- (z5);
\draw[arrow, zone5color] (z5) -- (d5);

% Feedback loops
\draw[arrow, dashed, red, bend right=30] (d5) to node[above, sloped, text=red] {Feedback} (z1);
\draw[arrow, dashed, red, bend right=20] (z3) to node[above, sloped, text=red] {Retraining} (z2);
\draw[arrow, dashed, red, bend left=20] (z5) to node[below, sloped, text=red] {Monitoring} (z4);

% Validation gates
\node[draw, circle, fill=green!30, minimum size=0.5cm] at ($(z1)!0.5!(z2)$) {\checkmark};
\node[draw, circle, fill=green!30, minimum size=0.5cm] at ($(z2)!0.5!(z3)$) {\checkmark};
\node[draw, circle, fill=green!30, minimum size=0.5cm] at ($(z3)!0.5!(z4)$) {\checkmark};

\end{tikzpicture}
\caption{Flux de données complet à travers les 5 zones}
\label{fig:dataflow}
\end{figure}

\subsection{Points de validation}

À chaque transition inter-zones, des \textbf{validation gates} assurent la conformité aux critères de sécurité :

\begin{enumerate}
    \item \textbf{Zone 1 $\rightarrow$ Zone 2}: Qualité données $\geq 70\%$, Pas de poisoning détecté
    \item \textbf{Zone 2 $\rightarrow$ Zone 3}: Entraînement convergé, Privacy budget respecté
    \item \textbf{Zone 3 $\rightarrow$ Zone 4}: Certification Grade $\geq$ A, Tous tests passés
    \item \textbf{Zone 4 $\rightarrow$ Zone 5}: Monitoring actif, Logs fonctionnels
\end{enumerate}

\newpage

\section{Tableau comparatif des zones}

\begin{table}[h!]
\centering
\caption{Synthèse comparative des 5 zones de sécurité}
\label{tab:zones_comparison}
\scriptsize
\begin{tabularx}{\textwidth}{|l|X|X|X|}
\hline
\textbf{Zone} & \textbf{Objectif Principal} & \textbf{Technologies Clés} & \textbf{Output} \\ \hline
\rowcolor{zone1color!10}
\textbf{Zone 1} & Garantir intégrité et confidentialité des données & SHA-256, Isolation Forest, AES-256, Tests statistiques & Dataset propre et validé \\ \hline
\rowcolor{zone2color!10}
\textbf{Zone 2} & Entraîner modèle robuste aux attaques adversariales & TRADES, Differential Privacy (Opacus), PGD, Docker & Modèle robuste entraîné \\ \hline
\rowcolor{zone3color!10}
\textbf{Zone 3} & Certifier robustesse et équité du modèle & FGSM, PGD, AutoAttack, Fairness metrics, Pytest & Modèle certifié Grade A+ \\ \hline
\rowcolor{zone4color!10}
\textbf{Zone 4} & Déployer en production avec défenses actives & FastAPI, Docker, WAF, JWT, Adversarial Detector & Service production sécurisé \\ \hline
\rowcolor{zone5color!10}
\textbf{Zone 5} & Surveiller et maintenir conformité continue & Prometheus, Grafana, Audit logging, RGPD compliance & Insights et alertes \\ \hline
\end{tabularx}
\end{table}

\section{Avantages de l'architecture}

\begin{tcolorbox}[colback=green!5, colframe=green!50!black, title=\textbf{Points Forts de l'Architecture}]
\begin{enumerate}
    \item \textbf{Défense en profondeur}: Plusieurs couches de sécurité indépendantes
    \item \textbf{Modularité}: Chaque zone peut être mise à jour indépendamment
    \item \textbf{Traçabilité complète}: Audit trail de bout en bout
    \item \textbf{Résilience}: Capacité de rollback et de recovery à chaque niveau
    \item \textbf{Conformité by design}: RGPD et réglementations intégrées
    \item \textbf{Monitoring continu}: Détection proactive des anomalies
    \item \textbf{Validation rigoureuse}: Certification avant déploiement
    \item \textbf{Reproductibilité}: Containérisation et versioning complets
\end{enumerate}
\end{tcolorbox}

\section{Limitations et travaux futurs}

\begin{tcolorbox}[colback=orange!5, colframe=orange!50!black, title=\textbf{Limitations Identifiées}]
\begin{itemize}
    \item \textbf{Overhead de performance}: Latence accrue de $\sim$30\% vs baseline
    \item \textbf{Complexité opérationnelle}: Nécessite expertise DevSecOps
    \item \textbf{Coût computationnel}: Adversarial training $\sim$2-3x plus long
    \item \textbf{Trade-off précision/robustesse}: Légère perte d'accuracy ($\sim$2-3\%)
\end{itemize}
\end{tcolorbox}

\begin{tcolorbox}[colback=blue!5, colframe=blue!50!black, title=\textbf{Extensions Possibles}]
\begin{itemize}
    \item \textbf{Federated Learning}: Entraînement distribué sans centralisation
    \item \textbf{Explainability}: Intégration SHAP/LIME pour interprétabilité
    \item \textbf{Attaques adaptatives}: Défenses contre adversaires conscients
    \item \textbf{Quantum resistance}: Préparation cryptographie post-quantique
    \item \textbf{Automated retraining}: Pipeline MLOps avec retraining automatique
\end{itemize}
\end{tcolorbox}

\section{Conclusion}

L'architecture de sécurisation à 5 zones proposée offre un framework complet et modulaire pour le déploiement sécurisé de systèmes d'IA critiques. En intégrant des mécanismes de défense à chaque étape du cycle de vie ML, elle garantit :

\begin{itemize}
    \item Une \textbf{robustesse adversariale} prouvée ($>50\%$ d'amélioration)
    \item Une \textbf{protection de la confidentialité} via Differential Privacy
    \item Une \textbf{conformité réglementaire} (RGPD) by design
    \item Une \textbf{traçabilité complète} pour audit et forensics
    \item Une \textbf{résilience opérationnelle} en production
\end{itemize}

Cette architecture a été validée empiriquement sur un système de détection d'armes pour les forces de l'ordre, démontrant sa faisabilité technique et économique avec des services cloud gratuits.

\vfill

\begin{center}
\rule{0.8\textwidth}{0.5pt}\\[0.5em]
\textit{Document généré pour le mémoire}\\
\textit{``Sécurisation des Systèmes d'IA pour le Maintien de l'Ordre''}
\end{center}

\end{document}
